\startslide
\commonframe
\mbox{}

\begin{tikzpicture}[remember picture,overlay]
  \slidetitle{Итоговый вывод}

  \node[
    anchor=north west,
    fill=softbg,
    rounded corners=10pt,
    text width=13.7cm,
    inner xsep=0.35cm,
    inner ysep=0.30cm,
    align=left
  ] at ([xshift=1.00cm,yshift=-1.65cm]current page.north west) {%
    {\fontsize{10.5}{12.4}\selectfont
    Главное различие между рассматриваемыми решениями заключается не только в
    выборе СУБД, но и в выборе \textbf{модели хранения одного и того же журнала операций}.}%
  };

  \node[anchor=north west, align=left, text width=13.8cm] at ([xshift=1.05cm,yshift=-3.62cm]current page.north west) {%
    {\fontsize{10.0}{11.8}\selectfont
    \textbf{PostgreSQL} --- лучший кандидат для глобального operation log и аналитических выборок.\\[0.12cm]
    \textbf{MongoDB} --- естественный выбор для объектно-центричных сценариев и хранения полной истории объекта в близкой к приложению форме.\\[0.12cm]
    \textbf{YTsaurus} --- перспективен для распределённого масштаба, но требует
    очень аккуратного проектирования ключа и представлений под конкретные запросы.}%
  };

  \node[
    anchor=north west,
    fill=softcyan,
    rounded corners=10pt,
    text width=13.7cm,
    inner xsep=0.35cm,
    inner ysep=0.30cm,
    align=left
  ] at ([xshift=1.00cm,yshift=-6.65cm]current page.north west) {%
    {\fontsize{10.0}{11.8}\selectfont
    Практический вывод работы: схему хранения для op-based CRDT нужно выбирать
    от \textbf{доминирующего паттерна чтения}, а не только от класса СУБД.}%
  };
\end{tikzpicture}

\footerband{8}
